\section*{Önsöz}
\indent\indent Bu çalışmayı yapma fikri ve isteği, Leman Dergisi'nin 26 Haziran 2025 tarihli sayısında yer alan bir karikatürden doğdu. Karikatürde kanatlı iki figür(figürlerden biri Hz.Muhammed, diğeri Hz.Musa)
birbirine selam vererek el sıkışırken, arka alanda yeryüzüne bombalar yağması resmedilmişti. Karikatür açık bir şekilde o sırada devam etmekte olan İran-İsrail Savaşı'na atıfta bulunyordu. Karikatür, toplumda
bir infiale sebep olmuş ve dergi binası saldırıya uğramıştı. Bunun üzerine dergiye soruşturma açılmış ve soruşturma kapsamında karikatürün çizeri Doğan Pehlevan, derginin yazı işleri müdürleri Aslan Özdemir
ve Zafer Aknar, grafiker Cebrail Okçu ve müessese müdürü Ali Yavuz tutuklanmıştı.

Dergi çalışanları , 5237 Sayılı Türk Ceza Kanunu'nun 216/3 maddesi uyarınca tutuklanmıştı. Madde şöyledir: \textit{Halkın bir kesiminin benimsediği dini değerleri alenen aşağılayan kişi, fiilin kamu barışını bozmaya elverişli olması halinde, altı aydan bir yıla kadar hapis cezası ile cezalandırılır.}
Görüldüğü üzere madde, halkın bir kesiminin benimsediği dini değerleri alenen aşağılayan kişiyi cezalandırmayı amaçlamaktadır. Hz.Muhammed'i bir figür olarak resmetmek, İslam dini değerlerini alenen aşağılamak mıdır? 
Yahut hangi eylemler alenen aşağılama kapsamına girmektedir? Aşağılama ile alenen aşağılama arasındaki fark nedir? Burada karar, hakimin yorumuna mı bırakılmıştır? Eğer öyleyse, hakimin yorum yaptığı bir sistemde yasalar 
nasıl herkes için eşit şekilde uygulanabilir? 

Bu soruların cevabını bulmak amacıyla Türkiye'de kazuistik hukuk kavramının nasıl ve ne şartlarda ortaya çıktığını araştırmak gerektiğini düşündüm. Şer`î mahkemeler, net bir kanun kodifikasyonunun olmadığı ve kadının içtihadı üzerinden yürüyen mahkemelerdi. Günümüzde ise ceza kanunu 
dediğimiz bir kitabın ihtiva ettiği suçlar ve cezalar üzerinden yargılama yapılmaktadır. Bu geçiş tam olarak nasıl olmuştu? Bunu anlamak amacıyla Osmanlı Devleti'nde kaleme alınmış ceza kanunlarının zamanla nasıl değiştiğini incelemeye karar verdim.

Çalışmama başlarken bana destek olan ve yol gösteren değerli hocam Prof. Dr. Arzu Terzi'ye ve dersin asistanı Aslı Tuna'ya teşekkür ederim. Bununla birlikte Leman Dergisi çalışanları Doğan Pehlevan, Aslan Özdemir, Zafer Aknar, Cebrail Okçu ve Ali Yavuz'a da verdikleri ilhamdan ötürü teşekkür ederim.

\clearpage